\documentclass[11pt]{article}
\usepackage{fontspec}
\setmainfont{TeX Gyre Pagella}
\usepackage[margin=1.3in]{geometry}
\usepackage{amsmath,amssymb}
\usepackage{microtype}
\usepackage[hidelinks]{hyperref}

\title{Unfinishable Games: Real-Time Structure as a Distinct Mechanism of Extraction}
\author{Flyxion \\ \small Independent Researcher}
\date{August 2026}

\begin{document}
\maketitle

\begin{abstract}
Criticism of exploitative game design has concentrated almost entirely on monetization, and above all on loot boxes, whose probability-weighted purchases have been analyzed as structurally close to gambling. This essay proposes a second, analytically distinct construct, structural unfinishability, present in a recognizable class of real-time, globally synchronous, alliance-based multiplayer games, and argues that it extracts a different resource through a different structural route than monetization does. A game is structurally unfinishable for a participant when consequential state changes can occur during that participant's absence, when the timing of those changes is substantially controlled by other players or the server rather than by the participant, when the resulting losses are difficult to reverse, and when responsibility for preventing those losses is socially attributed to that participant. The essay formalizes this as an exposure integral over a disengagement-cost function, distinguishes it from ordinary persistent-world design, and argues that the labor this structure elicits from coordinating players becomes part of what publishers sell as retention and engagement, an economic claim requiring its own defense rather than an assumed one. \emph{Last Shelter: Survival} is examined as a bounded case study illustrating the architecture rather than as proof that every synchronous game produces identical effects. The essay closes by distinguishing several possible design remedies and their respective enforcement difficulties, and by arguing that monetary and temporal extraction require different interventions, one addressed by probability disclosure and purchase regulation, the other addressed only by redesigning the temporal and social architecture of the game itself.
\end{abstract}

\section{Two Critiques, One Vocabulary}

The last decade of criticism directed at exploitative game design has converged almost entirely on a single target: monetization systems that use randomized rewards to convert player spending into a variable-ratio reinforcement schedule structurally similar to a slot machine. This body of criticism is substantial. Loot box mechanics have been associated with problem-gambling behavior \cite{zendle2019}, and consumer-protection bodies across multiple countries have characterized their sale and presentation as exploitative, particularly where the games in question are accessible to minors \cite{nccreport2022}. The international legal picture, however, is considerably more fragmented than a simple narrative of bans would suggest, a point developed in the next section, and scholars have documented substantial inconsistency in how different jurisdictions classify loot boxes as gambling or decline to \cite{xiao2021}.

What has received comparatively little independent treatment, at least outside the game-studies literature on unpaid guild labor discussed below, is a different mechanism, present in a recognizable class of games that happen to overlap with, but are not defined by, aggressive monetization. These are real-time, globally synchronous, alliance-based multiplayer titles, city-building and territory-control strategy games chief among them, in which meaningful in-game events can occur continuously, in real time, across a player base distributed over many time zones. Such a game can be entirely free to play and still impose severe, sustained coordination burdens on the people who invest themselves in it, because the burden in question is not a function of what the game charges. It is a function of how much control a participant retains over when it is safe to be absent.

Treating both mechanisms under a single heading, predatory monetization, or more loosely, addictive design, has a real cost. It leaves the temporal mechanism under-examined as its own analytic object, and it leaves any proposed remedy misaimed. Loot box regulation, purchase-probability disclosure, spending caps, and age gating are all responses calibrated to a gambling-adjacent purchase mechanic. None of them touches a game's real-time synchronous structure, and a game could in principle eliminate every randomized purchase it offers while leaving whatever temporal architecture it has completely intact.

\section{The Loot Box Model, Briefly}

The loot box model targets a specific, bounded resource, money, through a specific, well-characterized psychological mechanism. A player is offered the chance to obtain a desired item through a randomized draw rather than a fixed transaction, and the intermittent, variable schedule of reward this produces is among the most robust reinforcement patterns known to behavioral psychology \cite{king2018}. The harm is real and has been linked, at the population level, to problem-gambling scores \cite{zendle2019}, though the direction and mechanism of that association remains an active empirical question rather than a settled causal finding.

The regulatory response has been substantial but uneven rather than uniformly prohibitive. Belgium's gaming commission found certain paid loot box implementations to constitute gambling under national law. The Netherlands' gambling regulator issued a comparable finding against a specific implementation, but this was later overturned on administrative appeal, with the country's highest administrative court holding in 2022 that the loot box mechanic at issue did not constitute a standalone game of chance separable from the larger game \cite{xiao2021}. The United Kingdom's government-commissioned review concluded that an association between loot box spending and gambling-related harm existed but that a causal relationship had not been established, and it opted for industry-led protective measures rather than extending gambling law to cover the mechanic. A coalition of more than twenty consumer-protection organizations across eighteen European countries has nonetheless continued to call for stronger, harmonized regulation, characterizing existing loot box design as predatory and exploitative of cognitive biases \cite{nccreport2022}. The overall picture, in other words, is one of contested and evolving international treatment rather than a settled prohibition, and this essay does not take a position on how that contest should resolve.

None of this requires the game to run in real time. A turn-based game, or one with no live multiplayer component at all, can implement an identical loot box system with identical monetary extraction and identical gambling-adjacent psychology. The mechanism is orthogonal to the game's temporal structure. This orthogonality is the essay's starting point, because it means a second mechanism, tied specifically to real-time synchronous structure, cannot be a mere variant or intensification of the first. It has to be examined on its own terms, with its own formal apparatus, rather than folded into the borrowed vocabulary of addiction, which tends to relocate the explanatory burden onto the player's self-control rather than onto the properties of the system that make stopping costly.

\section{Formalizing Structural Unfinishability}

A precise definition needs to separate four phenomena that an informal account of ``real-time games'' tends to run together: persistent simulation, globally distributed participation, consequential events occurring during a participant's absence, and social responsibility for preventing collective losses. No one of these alone produces the mechanism this essay is concerned with. A persistent world can permit harmless absence. A globally distributed player base can coordinate asynchronously. Consequential events can be bounded by protection mechanics. Responsibility can be spread widely enough that no single participant remains continuously accountable. The mechanism this essay names arises specifically from their conjunction.

\begin{center}
\fbox{\begin{minipage}{0.9\textwidth}
A game is structurally unfinishable for participant $i$ when consequential state changes can occur during $i$'s absence, when their timing is substantially controlled by other players or the server rather than by $i$, when the resulting losses are difficult to reverse, and when responsibility for preventing those losses is socially attributed to $i$.
\end{minipage}}
\end{center}

This definition has two immediate consequences worth making explicit. First, it does not treat every persistent multiplayer game as equally unfinishable; a game satisfying only some of the four conditions, for instance one whose losses are always fully reversible, or whose responsibility is genuinely distributed across a large redundant group, falls outside the category even if it runs continuously and globally. Second, it makes unfinishability relational rather than a fixed property of the software: the same game can be an easily stoppable pastime for an ordinary member and an effectively continuous obligation for an alliance leader, because $R_i(t)$, the responsibility attributed to participant $i$, and $U_i(t)$, the difficulty of substituting another participant for $i$, both vary by role.

These quantities admit a formal treatment. A first pass might let $E(t) \in \{0,1\}$ indicate whether a consequential event can occur at time $t$, but a binary indicator conflates two games with very different actual exposure: one where an attack becomes possible once a month and one where it becomes possible every twenty minutes both register as $E=1$ almost continuously. The more useful quantity is a hazard or event-intensity term $\lambda_i(t)$, the expected rate of consequential events conditional on $i$'s absence at time $t$, which is estimable in principle from server logs rather than only from the binary fact that some window is open. Let $L_i(t)$ denote the expected loss to participant $i$ conditional on such an event occurring while $i$ is absent, let $R_i(t)$ denote the degree of responsibility socially attributed to $i$ for preventing that loss, and let $U_i(t)$ denote the difficulty of substituting another participant for $i$ at that moment. Define $i$'s temporal exposure over an interval $T$ as
\[
X_i(T) \;=\; \frac{1}{|T|} \int_T \lambda_i(t)\,L_i(t)\,R_i(t)\,U_i(t) \; dt.
\]
A safe disengagement interval $I \subseteq T$ exists when
\[
\int_I \lambda_i(t)\,L_i(t)\,R_i(t)\,U_i(t) \; dt \;\le\; \varepsilon
\]
for some tolerable threshold $\varepsilon$. This decomposition is useful beyond notational tidiness: it separates four independent levers, how exposed a participant is ($\lambda$), how costly an event would be ($L$), how responsible they are held ($R$), and how irreplaceable they are ($U$), and a design intervention can reduce unfinishability by attacking any one of the four without needing to touch the others, a correspondence taken up directly in the remedies developed later in the essay. On this formulation, structural unfinishability is not the literal absence of any interval of low exposure; it is the scarcity, or the participant's inability to unilaterally guarantee the existence, of a sufficiently long interval satisfying the inequality above. This makes the concept comparative rather than categorical and, in principle, measurable: a genre, a specific game, or a specific role within a game could be assessed by whether such intervals reliably exist and whether the participant, rather than the rest of the world's activity, controls where they fall.

This formal apparatus also draws a necessary line between persistence and unfinishability. A persistent world, one that continues to simulate and change while a given player is offline, is not automatically extractive in this essay's sense. What matters is whether $\lambda_i(t)L_i(t)R_i(t)U_i(t)$ is bounded away from zero for a given participant, not whether the world itself keeps running. A voxel-building server can operate continuously all night without $L_i(t)$ or $R_i(t)$ rising appreciably for any particular player, because nothing consequential and attributable to that player is at stake during those hours. Unfinishability names the specific condition in which the world's continuous operation is coupled to a specific participant's exposure rather than decoupled from it.

\section{The Circadian Vulnerability}

A further asymmetry explains why global synchronicity in particular, rather than mere real-time operation, aggravates $\lambda_i(t)$ for a globally distributed alliance. Human availability is approximately periodic: players sleep, and they sleep on a roughly twenty-four hour cycle tied to their local time zone. This periodicity, by itself, does not depend on how many adversaries a defender faces or where those adversaries are located; what changes with adversarial dispersion is the aggregate availability of the adversary population as a whole. Let $a_j(t)$ denote attacker $j$'s availability at time $t$, and let
\[
A(t) \;=\; \sum_j a_j(t)
\]
denote the aggregate availability of the full adversarial population. When attackers are concentrated in a single time zone, $A(t)$ inherits a strong daily rhythm, rising and falling together with that zone's local day. As the population's time-zone dispersion increases, the individual peaks and troughs of $a_j(t)$ increasingly fall at different points in the day, and the variance of the sum falls correspondingly,
\[
\operatorname{Var}_t[A(t)] \;\downarrow\; \text{as adversarial temporal dispersion increases,}
\]
even while each individual attacker retains their own strongly periodic availability. In the limit of a sufficiently large and well-distributed adversary population, $A(t)$ approaches a roughly constant level across the full twenty-four hour cycle. A defending player, meanwhile, retains their own strongly periodic availability, low during their local sleeping hours, high otherwise. The asymmetry this essay is pointing to is therefore not that any individual attacker specifically targets a defender's sleeping hours, but that the defender's periodic low-availability window is met by an adversarial population whose aggregate availability has been flattened by dispersion into something close to constant, so that the defender's least-available hours are, in expectation, no less contested than any other hours. Circadian absence, a biological constant no game design choice can remove, is thereby converted into a strategically exploitable regularity in $\lambda_i(t)$ purely as a function of the adversarial population's geographic dispersion, without requiring any single attacker to act with deliberate intent. This yields a testable prediction distinct from a generic account of engagement psychology: the severity of the effect should scale with the variance reduction in $A(t)$ as dispersion increases, a claim that a purely reinforcement-based account of engagement does not generate and that could in principle be checked against attack-timing data from specific games.

\section{Delegability}

Delegability, formalized above as low $U_i(t)$, determines whether global synchronicity actually forces continuous individual availability or can instead be absorbed by a group. If responsibility for responding to a threat can be handed seamlessly to another participant, perhaps one currently awake in a different time zone, then no single individual need remain continuously available even though $\lambda_i(t)$ itself never falls to zero. The configuration this essay is concerned with therefore requires a specific conjunction of high $\lambda_i$, high $L_i$, high $R_i$, and high $U_i$ simultaneously; removing any one of these, through protected rest periods, reliable delegation, bounded and recoverable losses, or a genuinely redundant responsibility structure, moves a game outside the category even if it remains globally synchronous. This is a stronger and more falsifiable claim than an unqualified assertion that real-time games extract time, because it predicts which superficially similar games should escape the mechanism, not only which ones exhibit it.

\section{From Session to Membership}

The social root of this structure deserves treatment as its own mechanism. A server clock alone cannot make a person wake at an unscheduled hour; what supplies the compulsion, formally, is $R_i(t)$, and $R_i(t)$ is a social quantity, assigned by other participants who are waiting, depending on the player's presence, or standing to suffer a consequence alongside them. The clock and the adversary population determine when $\lambda_i(t)$ is high; the social structure determines whether $R_i(t)$ attaches to a particular person once it is.

This is clearest by contrast with an older model of multiplayer participation, exemplified by real-time strategy titles played over a modem or local network, in which multiplayer was an optional session: players deliberately connected, played a bounded match, and disconnected, with no persistent structure surviving the session's end. A modern alliance-based persistent game replaces the match with membership. The player does not connect for a bounded engagement; they join a standing collective that continues to exist, accrue obligations, and face threats whether or not the player is currently present.
\[
\boxed{\text{multiplayer as an optional session} \;\longrightarrow\; \text{persistent social membership.}}
\]
Structural unfinishability, in the fullest sense this essay intends, is what results when membership of this kind is combined with the exposure structure formalized above: the game has stopped being a session one opts into and become a social position one occupies continuously, with occupancy itself carrying consequences for others.

\section{Persistent Games as Social-Platform Infrastructure}

The transition from session to membership has an implication worth drawing out explicitly rather than leaving embedded in the previous section, because it bears directly on where $R_i(t)$ comes from in the first place. Once a game supplies persistent identity, standing group membership, user-to-user communication, reputation or rank, and continuing interpersonal relationships, the software is no longer performing only the functions ordinarily associated with a game. The claim this essay needs is considerably more modest than an identity claim of the form massively multiplayer game equals social medium; it is a conditional one.
\begin{center}
\fbox{\begin{minipage}{0.9\textwidth}
\centering
persistent identity + communication + social graph + standing group membership $\;\Longrightarrow\;$ social-platform functions
\end{minipage}}
\end{center}
This matters for the present argument because $R_i(t)$ cannot arise from a server clock alone. A clock can make $\lambda_i(t)$ nonzero at a given hour, but it cannot, by itself, attach responsibility for the consequences of that hour to any specific person. That attribution requires exactly the durable relational infrastructure named above, an alliance that remembers who holds which rank, a chat channel through which expectations are communicated and enforced, a reputation system through which absence becomes visible to others. The platform layer is what manufactures the durable relationships through which responsibility becomes attributable to $i$ in particular rather than diffusing across an anonymous population.

This also sharpens the right-to-absence criterion developed later in the essay, by connecting it to a concern already familiar from criticism of social platforms rather than treating it as an idiosyncratic complaint about one game genre. Social platforms have long been criticized for making a user's absence itself socially legible and costly, missed notifications, missed conversations, a sense of falling behind a circulating conversation. The architecture examined in this essay compounds that familiar mechanism with a second, mechanical one that ordinary social media lacks: not only can a participant feel that they are missing something during an absence, their assets, territory, or alliance's competitive standing can be materially changed by other participants specifically because they were absent. The social consequence of absence and the mechanical consequence of absence reinforce rather than merely coexist with one another, which is a stronger combination than either social platforms or older, non-persistent multiplayer games produce on their own.

\section{Extraction Requires a Beneficiary}

The argument so far establishes a burden, sustained temporal exposure and the coordination labor it elicits from certain roles, but burden alone does not establish extraction in any economically meaningful sense; extraction requires showing that the burden's product is captured by someone other than the person bearing it. This distinction is worth stating as its own boxed claim, because it is easy to let the word extraction do more work than the argument has actually earned.
\begin{center}
\fbox{\begin{minipage}{0.85\textwidth}
\centering
unfinishability $\;\neq\;$ extraction
\end{minipage}}
\end{center}
Unfinishability is the structural condition analyzed in the preceding sections; extraction occurs only when some actor captures value produced under that condition. The two are logically separable, and a counterexample makes the separation concrete: a nonprofit, open-source persistent game could in principle exhibit exactly the same $\lambda_i$, $L_i$, $R_i$, and $U_i$ profile, the same structural unfinishability, without any publisher capturing economic surplus from the resulting coordination labor, because no publisher exists to capture it. Unfinishability would remain a real burden on the players who experience it, but calling it extraction in that case would be a category error.

The relevant claim for a commercially operated game is therefore a distinct, additional hypothesis rather than a restatement of the structural one. Let $V_{\mathrm{social}}$ denote the value of the player-produced organization that coordination labor generates, recruitment, scheduling, dispute resolution, rivalry, and the general cohesion of an active player base, and consider the hypothesis
\[
\frac{\partial \, \mathrm{Retention}}{\partial \, V_{\mathrm{social}}} \;>\; 0,
\]
that a publisher's retention, and by extension its monetizable engagement, rises with the value of this uncompensated player-produced organization. Stated this way, extraction becomes an empirically testable economic hypothesis about a specific game and publisher, rather than a term the argument's own vocabulary presupposes. The plausible causal chain runs as follows.
\begin{center}
\fbox{\begin{minipage}{0.9\textwidth}
\centering
player coordination labor $\;\longrightarrow\;$ persistent social content $\;\longrightarrow\;$ retention and monetizable engagement
\end{minipage}}
\end{center}
An alliance leader who recruits and retains members, schedules events, moderates disputes, trains newcomers, and manufactures the rivalries that make ongoing play compelling is producing exactly the kind of durable social content that keeps a persistent game populated and worth continuing to operate, and worth continuing to spend on, for everyone else in that alliance. This labor is not compensated by the publisher, and its product, a functioning, socially cohesive, competitively engaged player base, is plausibly a genuine input to the game's commercial success. Whether the hypothesis above actually holds for a given commercial title, and how large the effect is relative to other drivers of retention, is an empirical matter this essay does not settle; what the essay establishes is the more limited claim that appropriation, a publisher benefiting functionally from an active, well-organized population, is not automatically the same thing as extraction, deliberate capture of the specific economic surplus that uncompensated coordination labor generates, and that a commercially operated persistent game exhibiting structural unfinishability should be examined for the latter rather than assumed to practice it merely because the former is present.

\section{Case Study: \emph{Last Shelter: Survival}}

\emph{Last Shelter: Survival} provides a concrete, bounded illustration of this architecture; it is presented here as an example, not as evidence that every game in its genre produces identical effects. The developer's own listing describes the game as a real-time strategy title built around persistent base development, alliance formation, and continuous conflict between players worldwide \cite{lastshelter_playstore}, and the current storefront listing discloses loot boxes, in-app purchases, and persistent messaging and chat functionality, with an age rating of thirteen and older \cite{lastshelter_appstore}.

The game's architecture combines several of the elements formalized above in a single artifact. Its account and alliance structures are persistent rather than session-bound: a player's base and alliance membership remain situated in a shared, continuously operating world whether or not that player is currently active. Its competitive structure includes recurring multi-day events, among them a server-versus-server ``State vs State'' competition that a community-maintained guide describes as running on a roughly forty-eight-hour cycle, during which an alliance's collective standing depends on continuous, coordinated participation across the event's full duration \cite{lastshelter_swvs}, alongside recurring alliance-scale content such as cooperative defense events and territory-control competitions documented in player-produced guides \cite{lastshelter_appgamer}. The game does provide a partial, individually scoped mitigation in the form of a temporary defensive shield a player can activate to protect their own base from attack, but this protection is subject to a reactivation cooldown and, as the same community guide documents, does not by itself relieve the coordination requirement borne by alliance leadership during a live, alliance-wide event \cite{lastshelter_swvs}.

Applying the formal apparatus above, an ordinary alliance member in this game may have comparatively low $R_i(t)$ and low $U_i(t)$: their individual absence is substitutable, and responsibility for the alliance's overall standing is not specifically attributed to them. An alliance leader, diplomat, or event coordinator occupies a different position on both quantities. Responsibility for war declarations, territorial planning, defensive scheduling, and member coordination is commonly and visibly attributed to specific ranked roles within the alliance hierarchy, and, particularly in smaller or less redundantly staffed alliances, the difficulty of substituting another participant for a given coordinator during an active event may be high. Unfinishability, on this reading, is not a claim about the software in the abstract but a claim about a specific, identifiable role within it, exactly the relational property the formal definition above was built to capture.

A caveat about the evidentiary weight each source carries here is worth being explicit about. Official storefront listings are appropriate evidence for the game's advertised architecture, its real-time and alliance-based design, its communication and purchase systems, its disclosed mechanics, and its age classification; they establish $\lambda_i(t)$ and, more weakly, $L_i(t)$ reasonably well. Player-produced community guides are reasonable evidence that particular mechanics, such as the shield and its cooldown, exist and function as described, but they are considerably weaker evidence for exact claims about strategic significance. Neither source establishes $R_i(t)$ or $U_i(t)$ directly, since both are social quantities that would properly require observation of actual alliance communications, interviews with coordinating players, or survey data to measure. The more accurate statement of what this case study shows is therefore not that the full mechanism has been demonstrated, but that \emph{Last Shelter: Survival} instantiates the technical preconditions, persistent state, adversarially timed events, and a hierarchical responsibility structure, under which the proposed mechanism can arise, while whether particular players in particular roles actually cross a specified $X_i$ threshold remains an empirical question the case study does not itself answer.
\[
\text{architecture observed} \;\Longrightarrow\; \text{mechanism predicted} \;\Longrightarrow\; \text{player study tests the prediction.}
\]
Read this way, the case study generates a research program rather than substituting for one: a next step would be to measure $R_i(t)$ and $U_i(t)$ directly for identified coordination roles in this or a comparable game, rather than to treat their presence as already established by the architecture alone.

This case does not by itself establish that the game causes chronic sleep disruption or comparable harm at the population level; that would require player interviews, activity logs, or survey data specific to this genre, none of which this essay presents. What the case does establish is that the structural preconditions this essay has formalized, exposure that persists during absence, timing substantially outside a participant's control, difficult-to-reverse loss, and a hierarchical structure capable of attributing responsibility to specific coordinating roles, are jointly present in the advertised architecture of a widely played example of the genre, which is the more modest and more defensible claim the argument actually needs.

\section{Coordination as Structural Incentive, Not Established Population-Level Outcome}

This structural account converges with an older strand of game-studies research on the unpaid labor performed by guild and alliance leadership in earlier large-scale multiplayer games. Writing on massively multiplayer role-playing games, Levine documents that leading a large player collective routinely crosses from recreational play into unpaid, intensely time-consuming work \cite{levine2008}, a phenomenon related to, though not identical with, what K\"ucklich termed ``playbour'' in his analysis of unpaid labor in game modding specifically \cite{kucklich2005}; the analogy between modding labor and alliance coordination labor is useful but should be treated as analogical rather than as direct evidence about the genre this essay examines. Nardi's ethnographic account of World of Warcraft documents a comparable dynamic among raid leaders and guild officers, whose responsibilities extend beyond any single session into standing, ongoing coordination obligations \cite{nardi2010}, and Yee's survey-based research program on guild leadership similarly identifies the role as a source of disproportionate time burden and stress relative to ordinary play \cite{yee2006}.

The comparison to this earlier literature should be drawn carefully rather than overstated. World of Warcraft's world was itself persistent, and its guilds could be, and often were, geographically distributed; the earlier literature does not describe a game whose clock paused between sessions. What was comparatively bounded in that earlier context was more specifically the coordination requirement surrounding the game's highest-consequence collective content, which was commonly organized around socially negotiated session boundaries, scheduled raid nights with defined start and end times, and lockouts limiting how often a given encounter could matter. The distinction this essay needs is therefore not scheduled synchronous play against continuous play as such, but the following, narrower one: earlier raid-centered games were persistent, but much of their highest-coordination content could be scheduled by the participating group itself, whereas in globally synchronous territory games, consequential adversarial actions can instead be initiated by opponents during the defending group's absence, on a timetable the defenders do not control.

It is also worth being explicit about the limits of the available evidence. The playbour literature supports an analogy between guild leadership and unpaid labor; it does not, by itself, establish population-level empirical claims about how commonly severe or chronic consequences occur among players of the specific genre this essay examines. This essay's claim is structural rather than epidemiological: the architecture described can impose unpaid, job-like coordination duties and may encourage sleep-disrupting patterns of vigilance, especially among players occupying leadership or defense-critical roles. Whether, and how often, that structural incentive actually produces chronic sleep disruption or comparable harm at the population level is a separate empirical question, requiring dedicated study of this genre specifically, that this essay does not claim to have answered.

\section{Remedies and Their Enforcement Difficulties}

The formal apparatus developed above yields a design criterion in addition to a diagnosis. Let player $i$ choose, at some time $t_0$ of their own selection, an absence interval of duration $\Delta$, denoted $\mathcal{A}_i(t_0, \Delta)$. A system provides a genuine right to absence if, subject only to reasonable frequency and duration limits, participant $i$ can invoke $\mathcal{A}_i$ unilaterally and thereby obtain
\[
X_i\bigl([t_0, t_0+\Delta]\bigr) \;\le\; \varepsilon.
\]
\begin{center}
\fbox{\begin{minipage}{0.9\textwidth}
A well-designed persistent game should let a participant unilaterally invoke a bounded absence interval $\mathcal{A}_i(t_0,\Delta)$ over which their temporal exposure remains below a tolerable threshold.
\end{minipage}}
\end{center}
The word doing the essential work in this criterion is unilaterally. If invoking the interval requires permission from an alliance leader, negotiation with an adversary, locating and briefing a substitute, or merely hoping that no one attacks during the window, the system has not supplied a right to absence; other participants are supplying protection contingently, which is a different and considerably weaker thing. This yields a simple diagnostic question that can be asked of any persistent game or specific role within one: can the participant press leave and have the consequences stop? For a modem-era match, the answer is essentially yes, the match and its stakes end when the connection does. For an ordinary session-based game today, the answer is typically still yes. For a persistent game with a genuinely unilateral protection mechanism, the answer can be yes as well. For the architecture this essay has examined, the answer is often no, or no with qualification, particularly for participants occupying the high-$R_i$, high-$U_i$ coordinating roles the case study identified.

This formalization is considerably stronger than the generic ``take a break'' prompts that many always-on applications already display after some hours of continuous use, since a break reminder that fires while the underlying game mechanics simultaneously threaten a participant's alliance with territorial loss for having heeded it does not satisfy the criterion above; the reminder is a suggestion layered on top of an architecture that still penalizes compliance, not an actual instance of $\mathcal{A}_i$.

Satisfying this criterion in practice is harder than it first appears, and the difficulty is worth stating rather than glossing over. Several distinct mechanisms could contribute, and each carries its own limitation. Individual protection windows, of the kind \emph{Last Shelter: Survival}'s shield mechanic already provides, reduce a single participant's personal exposure but do not, by themselves, reduce the collective responsibility a coordinating role still bears during an active event. Alliance-level vulnerability windows, in which an entire group's exposure is bounded rather than any one member's, address the collective case but raise the question of whose local time such a window should be anchored to when the alliance itself spans every time zone; a window chosen by majority vote may simply exclude a minority of members from ever receiving protection. Asynchronous attack declarations, in which an aggressor must announce intent before a delay period elapses, followed by a mutually available resolution window, could decouple the timing of consequential events from any single participant's sleep schedule, at some cost to the immediacy that gives adversarial play its stakes. Hard caps on the frequency or duration of consecutive high-stakes events limit cumulative exposure without eliminating any single event's intensity. Reversible or insured offline losses reduce $L_i(t)$ directly, weakening the exposure integral regardless of when an event occurs. Role-redundancy requirements, mandating that no single coordination role be occupied by only one substitutable person, would directly lower $U_i(t)$ for critical roles. None of these is a complete solution on its own, and a genre-wide assessment of which combination best restores a genuine right to absence, without eliminating the adversarial stakes that make the genre distinctive, would be a natural next step for design research rather than something this essay resolves.

The regulatory dimension is correspondingly less settled than the loot box case. Probability disclosure regulation concerns observable transactions with a clear point of sale; temporal extraction, as formalized here, has no equivalent transactional moment for a regulator to examine. Addressing it would more plausibly require product-design standards, an analogy to occupational scheduling protections, targeted child-protection rules given that many titles in this genre carry the same modest age ratings as loot-box-bearing games, or a disclosure requirement covering a game's expected availability demands on coordinating roles. These should be read as possible directions for further consideration rather than as an existing or clearly applicable legal category.

\section{Conclusion}

The critical and regulatory apparatus built around loot boxes has done real work, and nothing in this essay is intended to diminish it. But that apparatus was built to answer a specific question, whether a randomized purchase mechanic functions as disguised gambling, and it has no purchase on a separate question that real-time, globally synchronous, alliance-based multiplayer games raise independently: whether a game's temporal and social architecture, quite apart from anything it sells, can remove a participant's unilateral control over when disengagement is safe, and whether the resulting coordination labor becomes an uncompensated input to the game's commercial operation. The distinctive claim this essay has defended is not that such games consume a great deal of time, since an absorbing game someone chooses to play for many hours is not thereby exploitative, but that structural unfinishability transfers control over the boundaries of play from the participant to the platform and to the other participants inhabiting it, so that leaving becomes something that can be done to a person's standing rather than something the person can simply decide to do.
\[
\boxed{\text{Exploitability can arise from the topology of permissible disengagement.}}
\]
Stated this way, the phenomenon generalizes beyond games, to on-call work schedules, notification-driven messaging platforms, always-live financial markets, and social media systems that make absence itself visible and costly to a user's standing, without requiring that any of these systems be psychologically identical to a game or to one another. Money and time are extracted through different mechanisms, diagnosed by different evidence, and addressed by different remedies. Treating the loot box model as though it exhausted the category of exploitative design leaves this second mechanism, and the design criteria needed to correct it, invisible to exactly the critical vocabulary that would otherwise be positioned to name it.

There is, finally, a structural symmetry worth noting between the argument developed here and a companion argument made elsewhere about the handling of historical record. That argument holds that a system should not let a later correction retroactively rewrite an earlier commitment, that a past claim, once made, should be preserved even as further evidence constrains how it is interpreted. The present essay holds a mirror-image position about the future rather than the past: a persistent system should not let standing membership retroactively claim unlimited authority over a participant's future availability. One argument denies a system authority over what already happened; the other denies a system authority over when participation must happen next. Both reject the same underlying architectural move, a system quietly converting continuity, whether of a record or of a membership, into an unbounded claim, on the past in one case, on the future in the other.

\begin{thebibliography}{99}

\bibitem{king2018}
D.~L.~King and P.~H.~Delfabbro, ``Predatory Monetization Schemes in Video Games (e.g. `Loot Boxes') and Internet Gaming Disorder,'' \emph{Addiction}, vol.~113, no.~11, pp.~1967--1969, 2018. DOI: 10.1111/add.14286.

\bibitem{xiao2021}
L.~Y.~Xiao, ``Regulating Loot Boxes as Gambling? Towards a Combined Legal and Self-Regulatory Consumer Protection Approach,'' \emph{Interactive Entertainment Law Review}, vol.~4, no.~1, pp.~27--47, 2021. DOI: 10.4337/ielr.2021.01.02.

\bibitem{zendle2019}
D.~Zendle and P.~Cairns, ``Loot Boxes are Again Linked to Problem Gambling: Results of a Replication Study,'' \emph{PLOS ONE}, vol.~14, no.~3, 2019.

\bibitem{nccreport2022}
Norwegian Consumer Council (Forbrukerr\aa det), ``Insert Coin: How the Gaming Industry Exploits Consumers Using Loot Boxes,'' report backed by more than twenty consumer organizations across eighteen European countries, May 2022.

\bibitem{levine2008}
S.~Levine, ``The Full-Time Guild Master,'' \emph{Intersect}, vol.~1, no.~1, pp.~36--42, Stanford University, 2008.

\bibitem{kucklich2005}
J.~K\"ucklich, ``Precarious Playbour: Modders and the Digital Games Industry,'' \emph{Fibreculture Journal}, no.~5, 2005.

\bibitem{nardi2010}
B.~A.~Nardi, \emph{My Life as a Night Elf Priest: An Anthropological Account of World of Warcraft}, University of Michigan Press, 2010.

\bibitem{yee2006}
N.~Yee, ``The Daedalus Project: Life as a Guild Leader,'' \emph{The Daedalus Project: The Psychology of MMORPGs}, research archive, nickyee.com, 2006.

\bibitem{lastshelter_playstore}
Long Tech Network Limited, ``Last Shelter: Survival,'' Google Play Store listing, accessed 2026.

\bibitem{lastshelter_appstore}
IM30 Technology Limited, ``Last Shelter: Survival,'' Apple App Store listing, accessed 2026.

\bibitem{lastshelter_swvs}
``Last Shelter Survival: State vs State Guide,'' player-produced community guide, Steam Community, accessed 2026.

\bibitem{lastshelter_appgamer}
``Last Shelter: Survival Answers and Guides,'' player-produced FAQ documentation, AppGamer, accessed 2026.

\end{thebibliography}

\end{document}

